
\documentclass[11pt,a4paper]{article}

\usepackage[utf8]{inputenc}
\usepackage[T1]{fontenc}
\usepackage[french]{babel}
\usepackage[a4paper,margin=2.1cm]{geometry}
\usepackage{amsmath,amssymb,amsthm,mathtools}
\usepackage{lmodern}
\usepackage{xcolor}
\usepackage{hyperref}
\usepackage{enumitem}
\usepackage{fancyhdr}
\usepackage{titlesec}
\usepackage{tcolorbox}
\usepackage{multicol}

\hypersetup{
  colorlinks=true,
  linkcolor=black,
  urlcolor=blue
}

\setlength{\parindent}{0pt}
\setlength{\parskip}{0.5em}

\pagestyle{fancy}
\fancyhf{}
\fancyhead[L]{Fiche d'exercices --- Analyse fonctionnelle-PGB}
\fancyhead[R]{Terminale}
\fancyfoot[C]{\thepage}

\titleformat{\section}{\large\bfseries}{\thesection.}{0.5em}{}
\titleformat{\subsection}{\normalsize\bfseries}{\thesubsection.}{0.5em}{}

\newcounter{exo}
\renewcommand{\theexo}{\arabic{exo}}

\definecolor{appcolor}{RGB}{25,95,160}
\definecolor{geocolor}{RGB}{0,120,90}
\definecolor{demcolor}{RGB}{150,80,0}
\definecolor{prepcolor}{RGB}{130,30,120}

\newcommand{\stars}[1]{%
  \ifcase#1\or $\star☆☆☆☆$\or $\star\star☆☆☆$\or $\star\star\star☆☆$\or $\star\star\star\star☆$\or $\star\star\star\star\star$\fi
}

\newcommand{\tagAPP}{\textcolor{appcolor}{\textbf{APP}}}
\newcommand{\tagGEO}{\textcolor{geocolor}{\textbf{GEO}}}
\newcommand{\tagDEM}{\textcolor{demcolor}{\textbf{DEM}}}
\newcommand{\tagPREP}{\textcolor{prepcolor}{\textbf{PREP}}}

\newcommand{\exo}[3]{%
  \refstepcounter{exo}
  \begin{tcolorbox}[colback=white,colframe=black!70,boxrule=0.5pt,arc=1.5mm]
  \textbf{Exercice \theexo} \hfill [#1 \; \stars{#2}]\\
  \textbf{#3}
  \end{tcolorbox}
}

\newcommand{\R}{\mathbb{R}}

\begin{document}

\begin{center}
    {\LARGE \textbf{Fiche d'exercices --- Analyse fonctionnelle}}\\[0.4em]
    {\large Niveau Terminale, avec progression de difficulté}\\[0.6em]
    \rule{0.9\textwidth}{0.5pt}
\end{center}

\textbf{AIDE de lecture} type d'exo: APP=application, GEO= application géométrique, PREP= raisonnement plus poussé, DEM= démonstrations, les étoiles répresentent un gradient de difficulté de 1 à 5.

\section*{Exercices d'application directe}

\exo{\tagAPP}{2}{Domaines de définition}
Déterminer l'ensemble de définition des fonctions suivantes et leur dérivée si elle existe :
\begin{enumerate}[label=\alph*)]
    \item \(f(x)=\dfrac{\sqrt{x+2}}{x-1}\)
    \item \(g(x)=\ln(2x-3)\)
    \item \(h(x)=\sqrt{\dfrac{x-1}{x+2}}\)
    \item \(k(x)=\ln(4-x^2)\)
\end{enumerate}

\exo{\tagAPP}{1}{Images, antécédents, injectivité}
Pour chacune des fonctions suivantes :
\begin{enumerate}[label=\alph*)]
    \item calculer l'image de \(1\),
    \item déterminer les antécédents éventuels de \(3\),
    \item préciser si la fonction est injective sur son domaine.
\end{enumerate}
\[
f(x)=x^2+2x,\qquad g(x)=3x-1,\qquad h(x)=x^3-3.
\]

\exo{\tagAPP}{2}{Variations usuelles}
Donner les variations des fonctions suivantes sur leur domaine, puis justifier brièvement :
\begin{enumerate}[label=\alph*)]
    \item \(x\mapsto x^3-1\)
    \item \(x\mapsto \dfrac{1}{x^2+1}\)
    \item \(x\mapsto \sqrt{2x+3}\)
    \item \(x\mapsto \ln x\)
    \item \(x\mapsto e^{-x}\)
\end{enumerate}

\exo{\tagAPP}{2}{Dérivées}
Calculer les dérivées des fonctions suivantes :
\begin{enumerate}[label=\alph*)]
    \item \(f(x)=x^4-4x+1\)
    \item \(g(x)=\dfrac{x^2+1}{x}\)
    \item \(h(x)=e^{2x}+x^2e^x\)
    \item \(k(x)=\ln(x^2+3x+1)\)
    \item \(m(x)=\sqrt{3x-2}+\dfrac{1}{x}\)
\end{enumerate}

\exo{\tagAPP}{2}{Tangentes}
Déterminer l'équation de la tangente à la courbe de chacune des fonctions suivantes au point indiqué :
\begin{enumerate}[label=\alph*)]
    \item \(f(x)=x^2-2x+3\) en \(x=1\)
    \item \(g(x)=x^3-3x\) en \(x=1\)
    \item \(h(x)=e^x\) en \(x=1\)
    \item \(k(x)=\ln x\) en \(x=e\)
\end{enumerate}

\exo{\tagAPP}{2}{Étude simple d'un polynôme}
Étudier complètement la fonction
\[
f(x)=x^2-4x+\frac{11}{2}.
\]

\exo{\tagAPP}{2}{Fonction rationnelle}
Étudier la fonction
\[
f(x)=\frac{x^2+1}{x+1}.
\]

\exo{\tagAPP}{2}{Équations et inéquations}
Résoudre dans \(\R\), ou dans l'ensemble de définition adapté :
\begin{enumerate}[label=\alph*)]
    \item \(e^{2x+1}=3\)
    \item \(\ln(x-2)=1\)
    \item \(e^x\leq 5\)
    \item \(\ln(x+1)\geq 0\)
\end{enumerate}

\exo{\tagAPP}{2}{Fonctions composées}
Calculer les dérivées des fonctions suivantes :
\begin{enumerate}[label=\alph*)]
    \item \(f(x)=(x^2+1)^3\)
    \item \(g(x)=e^{x^2-1}\)
    \item \(h(x)=\ln(1+e^x)\)
    \item \(k(x)=\sqrt{x^2+4x+5}\)
\end{enumerate}

\exo{\tagAPP}{2}{Limites}
Calculer les limites suivantes :
\begin{enumerate}[label=\alph*)]
    \item \(\displaystyle\lim_{x\to+\infty}\frac{2x^2-3x+1}{x^2+1}\)
    \item \(\displaystyle\lim_{x\to+\infty}\bigl(x-\ln x\bigr)\)
    \item \(\displaystyle\lim_{x\to0^+}x\ln x\)
    \item \(\displaystyle\lim_{x\to+\infty}\frac{e^x}{x^3}\)
\end{enumerate}

\section*{Exercices d'analyse géométrique}

\exo{\tagGEO}{2}{Parabole}
Étudier la fonction
\[
f(x)=x^2-4
\]
et tracer soigneusement sa courbe.

\exo{\tagGEO}{2}{Asymptotes et branches}
Étudier la fonction
\[
f(x)=\frac{x+2}{x-1}
\]
et en tracer un croquis précis.

\exo{\tagGEO}{3}{Tangentes horizontales}
Étudier la fonction
\[
f(x)=x^3-3x+2
\]
et déterminer ses tangentes horizontales éventuelles.

\exo{\tagGEO}{3}{Convexité}
Étudier complètement la fonction
\[
f(x)=x^4-4x^2
\]
et en tracer un croquis précis.

\exo{\tagGEO}{3}{Logarithme composé}
Étudier complètement la fonction
\[
f(x)=\ln(1+x^2)
\]
et tracer sa courbe et préciser s'il y a des asymptotes particulières.

\exo{\tagGEO}{4}{Lecture graphique et paramètre}
Étudier la fonction et préciser s'il y a des asymptotes particulières.
\[
f(x)=x+\frac{4}{x}
\qquad \text{sur } ]0,+\infty[
\]
puis discuter, selon \(m\in\R\), le nombre de solutions de l'équation
\[
x+\frac{4}{x}=m.
\]

\section*{Exercices de démonstration}

\exo{\tagDEM}{2}{Monotonie par définition}
Montrer directement, à partir de la définition, que la fonction
\[
f(x)=5x-2
\]
est strictement croissante sur \(\R\).

\exo{\tagDEM}{2}{Fonction carré}
Montrer sans utiliser la dérivée :
\begin{enumerate}[label=\alph*)]
    \item que la fonction \(x\mapsto x^2\) est décroissante sur \(]-\infty,0]\),
    \item qu'elle est croissante sur \([0,+\infty[\),
    \item qu'elle n'est pas monotone sur \(\R\).
\end{enumerate}

\exo{\tagDEM}{2}{Monotonie et antécédents}
Montrer qu'une fonction strictement monotone sur un intervalle admet au plus un antécédent pour chaque réel.

\exo{\tagDEM}{3}{Racine carrée}
Montrer, sans utiliser la dérivée, que la fonction racine carrée est strictement croissante sur \([0,+\infty[\).

\exo{\tagDEM}{3}{Existence et unicité}
On considère la fonction
\[
f(x)=x^3+x-1.
\]
Montrer que l'équation \(f(x)=0\) admet une unique solution réelle, puis en donner un encadrement simple.

\exo{\tagDEM}{3}{Discussion selon un paramètre}
On considère la fonction
\[
f(x)=x^3-3x.
\]
Résoudre \(f(x)=0\), puis discuter, selon \(m\in\R\), le nombre de solutions de l'équation
\[
f(x)=m.
\]

\exo{\tagDEM}{3}{Raisonnement par l'absurde}
Montrer par l'absurde que l'équation
\[
e^x=-1
\]
n'admet aucune solution réelle.

\exo{\tagDEM}{3}{Inégalité}
Montrer que, pour tout réel \(x>0\),
\[
x-\ln x\geq 1.
\]
Déterminer le cas d'égalité.

\exo{\tagDEM}{4}{Suite récurrente}
On considère la suite \((u_n)\) définie par
\[
u_{n+1}=\frac{u_n+3}{2},
\qquad u_0=10.
\]
Étudier la convergence de cette suite et déterminer sa limite.

\exo{\tagDEM}{4}{Suite récurrente avec racine}
On considère la suite \((u_n)\) définie par
\[
u_{n+1}=\sqrt{u_n+2},
\qquad u_0=2.
\]
Étudier la convergence de cette suite.

\exo{\tagDEM}{4}{Suite de fonctions}
Pour tout entier \(n\geq 1\), on définit
\[
f_n(x)=\frac{x^{2n+1}}{1+x^{2n}}
\qquad \text{sur } \R.
\]
Pour tout réel \(x\), étudier la convergence de la suite \(\bigl(f_n(x)\bigr)\).  
On pourra également préciser l'allure générale de \(f_n\) pour un entier \(n\) fixé.

\exo{\tagDEM}{5}{Étude de racines}
On considère la fonction
\[
f(x)=e^x-x-2.
\]
Montrer que l'équation \(f(x)=0\) admet au moins une solution réelle, puis déterminer le nombre exact de solutions réelles.

\section*{Exercices de réflexion }

\exo{\tagPREP}{3}{Itération affine}
Déterminer toutes les fonctions affines \(f:\R\to\R\) de la forme
\[
f(x)=ax+b
\]
telles que
\[
f(f(x))=x
\qquad \text{pour tout } x\in\R.
\]

\exo{\tagPREP}{3}{Équation fonctionnelle}
Déterminer toutes les fonctions dérivables \(f:\R\to\R\) telles que
\[
f(x+y)=f(x)+f(y)
\qquad \text{pour tous } x,y\in\R.
\]

\exo{\tagPREP}{4}{Équation fonctionnelle enrichie}
Déterminer toutes les fonctions dérivables \(f:\R\to\R\) telles que
\[
f(x+y)=f(x)+f(y)+xy
\qquad \text{pour tous } x,y\in\R.
\]

\exo{\tagPREP}{4}{Exponentielle contre polynôme}
Étudier le nombre de solutions de l'équation
\[
x^2=e^x.
\]

\exo{\tagPREP}{4}{Équation paramétrée}
Soit \(m\in\R\). Discuter, selon \(m\), le nombre de solutions de l'équation
\[
x+\ln x=m
\qquad \text{sur } ]0,+\infty[.
\]

\exo{\tagPREP}{4}{Étude complète avec paramètre}
Étudier complètement la fonction
\[
f(x)=x-\ln x
\qquad \text{sur } ]0,+\infty[
\]
puis discuter, selon \(m\in\R\), le nombre de solutions de l'équation
\[
x-\ln x=m.
\]

\exo{\tagPREP}{4}{Fonction rationnelle}
Étudier complètement la fonction
\[
f(x)=\frac{x^2+2x+3}{x+1}
\]
et en tracer l'allure.

\exo{\tagPREP}{4}{Fonction bornée}
Étudier complètement la fonction
\[
f(x)=\frac{x^2}{x^2+1}
\]
et tracer sa courbe.

\exo{\tagPREP}{5}{Grande étude de fonction}
On considère la fonction
\[
f(x)=x^3-3x.
\]
Étudier complètement cette fonction, puis discuter, selon \(m\in\R\), le nombre de solutions de l'équation
\[
x^3-3x=m.
\]

\exo{\tagPREP}{5}{Fonction polynomiale paire}
Étudier complètement la fonction
\[
f(x)=x^4+x^2+1
\]
puis discuter, selon \(m\in\R\), le nombre de solutions de l'équation
\[
x^4+x^2+1=m.
\]

\exo{\tagPREP}{5}{Asymptote oblique et variations}
Étudier complètement la fonction
\[
f(x)=\frac{x^2+1}{x-1}
\]
et tracer un croquis précis de sa courbe.

\exo{\tagPREP}{5}{Synthèse}
On considère la fonction
\[
f(x)=e^x-x^2.
\]
Étudier complètement cette fonction, puis montrer que l'équation
\[
e^x=x^2
\]
admet exactement deux solutions réelles, que l'on encadrera.

\section*{Pour aller plus loin}

\exo{\tagPREP}{5}{Version plus ouverte}
Déterminer toutes les fonctions dérivables \(f:\R\to\R\) telles que
\[
f(x+y)-f(x)-f(y)
\]
ne dépende que du produit \(xy\). Proposer ensuite des exemples explicites.

\exo{\tagDEM}{4}{Convexité et tangente}
On considère la fonction
\[
f(x)=\ln(1+x^2).
\]
Déterminer la position de la courbe par rapport à sa tangente en \(0\), puis interpréter géométriquement le résultat.

\exo{\tagGEO}{3}{Lecture graphique comparée}
Comparer graphiquement les courbes des fonctions
\[
x\mapsto x^2,\qquad x\mapsto e^x,\qquad x\mapsto \ln(1+x^2)
\]
sur des intervalles pertinents. On mettra en évidence les comportements aux bornes, les croissances et les allures locales.

\exo{\tagAPP}{2}{Révisions ciblées}
Résoudre les équations suivantes :
\begin{enumerate}[label=\alph*)]
    \item \(\ln(x^2+1)=0\)
    \item \(e^x=x+1\)
    \item \(x-\dfrac{1}{x}=0\)
    \item \(\dfrac{x^2+1}{x-1}=2\)
\end{enumerate}

\exo{\tagDEM}{4}{Suite récurrente paramétrée}
Soit \(a>0\) et \((u_n)\) définie par
\[
u_{n+1}=\frac{u_n+a}{2}.
\]
Étudier le comportement de la suite selon la valeur de \(u_0\), puis déterminer sa limite éventuelle.

\exo{\tagPREP}{4}{Équation mixte}
Étudier le nombre de solutions de l'équation
\[
\ln x = 2-x
\qquad \text{sur } ]0,+\infty[.
\]

\exo{\tagPREP}{4}{Variation et extremum}
Étudier complètement la fonction
\[
f(x)=x+\frac{1}{x}
\qquad \text{sur } ]0,+\infty[
\]
puis résoudre, selon \(m\in\R\), l'équation
\[
x+\frac{1}{x}=m.
\]

\exo{\tagDEM}{3}{TVI}
Montrer que l'équation
\[
x^3-x-1=0
\]
admet au moins une solution dans l'intervalle \([1,2]\).

\exo{\tagDEM}{3}{TVI et unicité}
Montrer que l'équation
\[
x^3-x-1=0
\]
admet une unique solution réelle.

\end{document}
